\documentclass[12pt,a4paper]{article}
\usepackage[utf8]{inputenc}
\usepackage[russian]{babel}
\usepackage{amsmath, amssymb}
\usepackage{geometry}
\usepackage{booktabs}
\usepackage{array}
\usepackage{graphicx}
\usepackage{hyperref}
\usepackage{xcolor}
\usepackage{listings}
\usepackage{caption}
\usepackage{float}

\geometry{left=2.5cm, right=2cm, top=2cm, bottom=2cm}

\lstset{
  basicstyle=\small\ttfamily,
  backgroundcolor=\color{gray!10},
  frame=single,
  breaklines=true,
  inputencoding=utf8,
  extendedchars=true
}

\title{\textbf{Лабораторная работа №5}\\
\large Аппроксимация табличных данных:\\задачи оптимизации и машинного обучения\\[0.5em]
\normalsize Методы оптимизации, ИТМО, 2 курс, 4 семестр}
\author{%
  Вариант 9\\[0.4em]
  Студенты: Калинин Дмитрий, Рзаев Ахмед\\
  Преподаватель: Селина Е.\,Г.
}
\date{}

\begin{document}
\maketitle
\tableofcontents
\newpage

%=============================================================================
\section{Данные по варианту}
%=============================================================================

\begin{table}[H]
\centering
\caption{Датасет, вариант 9}
\begin{tabular}{cccc}
\toprule
$i$ & $x^i$ & $y^i$ & $z^i$ \\
\midrule
1 & 1.02 & 0.92 & 0.25 \\
2 & 2.15 & 1.88 & 3.02 \\
3 & 3.14 & 3.09 & 5.72 \\
4 & 4.13 & 4.07 & 1.51 \\
5 & 4.94 & 5.15 & 0.28 \\
\bottomrule
\end{tabular}
\end{table}

Задача: найти модельную функцию $z \approx z_{\text{model}}(x, y)$, наилучшим образом
аппроксимирующую табличные данные в смысле минимизации MSE.


%=============================================================================
\section{Задание 1. Аппроксимация двумерной гауссианой}
%=============================================================================

\subsection{Модель}

\begin{equation}
  z_G = A\,e^{-Q} + z_0, \qquad
  Q = \frac{\tilde{x}^2}{2\sigma_x^2} + \frac{\tilde{y}^2}{2\sigma_y^2},
\end{equation}
\begin{equation}
  \tilde{x} = (x - x_0)\cos\theta + (y - y_0)\sin\theta, \qquad
  \tilde{y} = -(x - x_0)\sin\theta + (y - y_0)\cos\theta.
\end{equation}

Параметры модели: $(A,\, x_0,\, y_0,\, \sigma_x,\, \sigma_y,\, \theta,\, z_0) \in \mathbb{R}^7$.

\subsection{Постановка задачи оптимизации}

\begin{equation}
  L = \frac{1}{2 \cdot 5}\sum_{i=1}^{5}\bigl(z_G^i - z^i\bigr)^2
  \;\longrightarrow\; \min_{A,\, x_0,\, y_0,\, \sigma_x,\, \sigma_y,\, \theta,\, z_0}.
\end{equation}

Метод: \textbf{L-BFGS-B}. Начальное приближение:
\[
  A_0 = z_{\max} + 0.1 = 5.82, \quad
  (x_0, y_0) = (3.14,\, 3.09), \quad
  \sigma_x = 0.70,\quad \sigma_y = 0.75, \quad
  \theta = 0, \quad z_0 = 0.
\]

\subsection{Результаты оптимизации}

\begin{table}[H]
\centering
\caption{Найденные параметры гауссианы}
\begin{tabular}{cc}
\toprule
Параметр & Значение \\
\midrule
$A$        & 5.9050 \\
$x_0$      & 2.9105 \\
$y_0$      & 2.5006 \\
$\sigma_x$ & 0.4546 \\
$\sigma_y$ & 1.7727 \\
$\theta$   & $-0.34$ рад $= -19.72°$ \\
$z_0$      & 0.1828 \\
Финальный MSE & 0.00000000 \\
\bottomrule
\end{tabular}
\end{table}

\subsection{Аналитический вид модели}

\begin{equation}
  z(x,y) = 5.9050\cdot\exp(-Q) + 0.1828,
\end{equation}
\begin{equation}
  Q = \frac{\tilde{x}^2}{2\cdot 0.4546^2} + \frac{\tilde{y}^2}{2\cdot 1.7727^2},
\end{equation}
\begin{align}
  \tilde{x} &= (x - 2.9105)\cos(-0.34) + (y - 2.5006)\sin(-0.34), \\
  \tilde{y} &= -(x - 2.9105)\sin(-0.34) + (y - 2.5006)\cos(-0.34).
\end{align}

\subsection{Невязки на каждом объекте}

\begin{table}[H]
\centering
\caption{Невязки гауссианы}
\begin{tabular}{ccccc}
\toprule
$i$ & $(x^i, y^i)$ & $z^i$ & $z_G^i$ & невязка \\
\midrule
1 & (1.02, 0.92) & 0.25 & 0.25 & 0.00 \\
2 & (2.15, 1.88) & 3.02 & 3.02 & 0.00 \\
3 & (3.14, 3.09) & 5.72 & 5.72 & 0.00 \\
4 & (4.13, 4.07) & 1.51 & 1.51 & 0.00 \\
5 & (4.94, 5.15) & 0.28 & 0.28 & 0.00 \\
\bottomrule
\end{tabular}
\end{table}

Невязки равны нулю (машинная точность) --- модель точно интерполирует все 5 точек.

\subsection{Визуализация}

\begin{figure}[H]
  \centering
  \includegraphics[width=\textwidth]{model_plot.png}
  \caption{Задание 1: кривая обучения, 3D-поверхность и линии уровня двумерной гауссианы}
\end{figure}


%=============================================================================
\section{Задание 2. Аппроксимация эллиптическим параболоидом}
%=============================================================================

\subsection{Модель}

Эллиптический параболоид в неканонической системе координат. Для варианта 9
экспериментальные точки имеют максимум в середине, поэтому используется параболоид,
направленный вниз:
\begin{equation}
  z_P = z_0 - a(x - x_0)^2 - b(y - y_0)^2 - 2c(x - x_0)(y - y_0),
  \quad ab > c^2.
\end{equation}

Параметры модели: $(x_0,\, y_0,\, z_0,\, a,\, b,\, c) \in \mathbb{R}^6$.

\subsection{Постановка задачи оптимизации}

\begin{equation}
  L = \frac{1}{2 \cdot 5}\sum_{i=1}^{5}\bigl(z_P^i - z^i\bigr)^2
  \;\longrightarrow\; \min_{x_0,\, y_0,\, z_0,\, a,\, b,\, c}.
\end{equation}

Метод: \textbf{L-BFGS-B}. Чтобы ограничение $ab > c^2$ выполнялось автоматически,
квадратичная форма параметризована через разложение Холецкого:
\[
  a = l_{11}^2 + \varepsilon,\qquad
  b = l_{21}^2 + l_{22}^2 + \varepsilon,\qquad
  c = l_{11}l_{21}.
\]
Начальное приближение: вершина рядом с точкой с максимальным $z^i$,
т.е.\ $(x_0, y_0, z_0) = (3.14,\, 3.09,\, 5.72)$.

\subsection{Результаты оптимизации}

\begin{table}[H]
\centering
\caption{Найденные параметры параболоида}
\begin{tabular}{cc}
\toprule
Параметр & Значение \\
\midrule
$x_0$ & 4.1529 \\
$y_0$ & 6.0076 \\
$z_0$ & 18.2989 \\
$a$   & 9.9725 \\
$b$   & 4.6181 \\
$c$   & $-6.2541$ \\
Проверка $ab > c^2$ & $46.0542 > 39.1132$ (True) \\
Финальный MSE & $5.4\cdot 10^{-11}$ \\
\bottomrule
\end{tabular}
\end{table}

\textbf{Замечание.} Вершина параболоида оказалась вне прямоугольника $[0,6]\times[0,6]$.
Это нормально для аппроксимации по пяти точкам: поверхность проходит через объекты
с машинной точностью, а ограничение эллиптичности выполняется.

\subsection{Аналитический вид модели}

\begin{equation}
  z(x,y) = 18.2989 - 9.9725\,(x - 4.1529)^2 - 4.6181\,(y - 6.0076)^2
            - 2(-6.2541)\,(x - 4.1529)(y - 6.0076).
\end{equation}

\subsection{Невязки на каждом объекте}

\begin{table}[H]
\centering
\caption{Невязки параболоида}
\begin{tabular}{ccccc}
\toprule
$i$ & $(x^i, y^i)$ & $z^i$ & $z_P^i$ & невязка \\
\midrule
1 & (1.02, 0.92) & 0.25 & 0.2500 & $+0.0000$ \\
2 & (2.15, 1.88) & 3.02 & 3.0200 & $-0.0000$ \\
3 & (3.14, 3.09) & 5.72 & 5.7200 & $+0.0000$ \\
4 & (4.13, 4.07) & 1.51 & 1.5100 & $+0.0000$ \\
5 & (4.94, 5.15) & 0.28 & 0.2800 & $-0.0000$ \\
\bottomrule
\end{tabular}
\end{table}

Невязки близки к нулю --- параболоид хорошо описывает данные, если разрешить
параболоид с максимумом, а не только чашу с минимумом.

\subsection{Визуализация}

\begin{figure}[H]
  \centering
  \includegraphics[width=\textwidth]{задание2_параболоид.png}
  \caption{Задание 2: кривая обучения, 3D-поверхность и линии уровня
  эллиптического параболоида.}
\end{figure}


%=============================================================================
\section{Задание 3. Константная модель}
%=============================================================================

\subsection{Минимизация MSE}

Лучшая константная модель в смысле минимизации MSE~--- \textbf{среднее арифметическое}:
\[
  \frac{d}{da}\,\frac{1}{2m}\sum_{i=1}^m (z^i - a)^2 = 0
  \;\Rightarrow\;
  a^* = \frac{1}{m}\sum_{i=1}^m z^i = \frac{0.25 + 3.02 + 5.72 + 1.51 + 0.28}{5} = \mathbf{2.1560}.
\]

\[
  \mathrm{MSE}(a^*) = 4.2036, \qquad
  \text{невязки: } [-1.906,\; 0.864,\; 3.564,\; -0.646,\; -1.876].
\]

\subsection{Минимизация MAE}

Лучшая константная модель в смысле минимизации MAE~--- \textbf{медиана}:
\[
  a^* = \mathrm{median}\{0.25,\; 0.28,\; 1.51,\; 3.02,\; 5.72\} = \mathbf{1.5100}.
\]

\[
  \mathrm{MAE}(a^*) = 1.6420, \qquad
  \text{невязки: } [1.26,\; 1.51,\; 4.21,\; 0.00,\; 1.23].
\]

\subsection{Сравнение и проверка}

\begin{table}[H]
\centering
\caption{Сравнение константных моделей}
\begin{tabular}{lcc}
\toprule
Критерий & Оптимальная константа & Значение Loss \\
\midrule
MSE & $a^* = 2.1560$ (среднее) & $\mathrm{MSE} = 4.2036$ \\
MAE & $a^* = 1.5100$ (медиана) & $\mathrm{MAE} = 1.6420$ \\
\bottomrule
\end{tabular}
\end{table}

Проверки:
\begin{itemize}
  \item $\mathrm{MSE}(\text{среднее}) = 4.2036 \leq \mathrm{MSE}(\text{медиана}) = 4.6209$ --- \textbf{True}
  \item $\mathrm{MAE}(\text{медиана}) = 1.6420 \leq \mathrm{MAE}(\text{среднее}) = 1.7712$ --- \textbf{True}
\end{itemize}


%=============================================================================
\section{Задание 4. Ручные вычисления: RBF-сеть (Стратегия 2)}
%=============================================================================

\subsection{Модель RBF-сети}

\begin{equation}
  z_{\mathrm{RBF}} = w_1\exp\!\left(-\frac{(x-c_{1x})^2+(y-c_{1y})^2}{2\sigma_1^2}\right)
  + w_2\exp\!\left(-\frac{(x-c_{2x})^2+(y-c_{2y})^2}{2\sigma_2^2}\right) + w_0.
\end{equation}

\subsection{Шаг 1. K-Means (k=2): поиск центров}

Начальные центры (случайный выбор): $c_1^{(0)} = (1.02,\;0.92)$, $c_2^{(0)} = (3.14,\;3.09)$.

\begin{table}[H]
\centering
\caption{Итерация it=0: назначение кластеров}
\begin{tabular}{ccccc}
\toprule
$i$ & $(x^i, y^i)$ & $\|p^i - c_1\|$ & $\|p^i - c_2\|$ & Кластер \\
\midrule
1 & (1.02, 0.92) & \textbf{0.0000} & 3.0337 & $C_1$ \\
2 & (2.15, 1.88) & \textbf{1.4827} & 1.5634 & $C_1$ \\
3 & (3.14, 3.09) & 3.0337 & \textbf{0.0000} & $C_2$ \\
4 & (4.13, 4.07) & 4.4266 & \textbf{1.3930} & $C_2$ \\
5 & (4.94, 5.15) & 5.7671 & \textbf{2.7356} & $C_2$ \\
\bottomrule
\end{tabular}
\end{table}

Обновление центров:
\[
  c_1^{(1)} = \frac{(1.02,0.92)+(2.15,1.88)}{2} = (1.585,\;1.400),
\]
\[
  c_2^{(1)} = \frac{(3.14,3.09)+(4.13,4.07)+(4.94,5.15)}{3} = (4.070,\;4.103).
\]

\begin{table}[H]
\centering
\caption{Итерация it=1: проверка сходимости}
\begin{tabular}{ccccc}
\toprule
$i$ & $(x^i, y^i)$ & $\|p^i - c_1\|$ & $\|p^i - c_2\|$ & Кластер \\
\midrule
1 & (1.02, 0.92) & \textbf{0.7414} & 4.4086 & $C_1$ \\
2 & (2.15, 1.88) & \textbf{0.7414} & 2.9376 & $C_1$ \\
3 & (3.14, 3.09) & 2.2965 & \textbf{1.3754} & $C_2$ \\
4 & (4.13, 4.07) & 3.6886 & \textbf{0.0686} & $C_2$ \\
5 & (4.94, 5.15) & 5.0318 & \textbf{1.3610} & $C_2$ \\
\bottomrule
\end{tabular}
\end{table}

Состав кластеров не изменился $\Rightarrow$ \textbf{центры найдены}:
\[
  \boxed{c_1 = (1.585,\;1.400), \quad c_2 = (4.070,\;4.103).}
\]

\subsection{Шаг 2. Начальное приближение}

\[
  |c_1 - c_2| = \sqrt{(4.070-1.585)^2 + (4.103-1.400)^2} = 3.6720,
\]
\[
  \sigma_1^{(0)} = \sigma_2^{(0)} = \frac{3.6720}{2} = 1.8360,
\]
\[
  w^{(0)} = (w_0, w_1, w_2) = (0.0,\;0.1,\;-0.1), \quad \eta = 0.2.
\]

\subsection{Шаг 3. Первая итерация градиентного спуска}

Значения базисных функций $G_j^i = \exp\!\left(-\dfrac{(x^i-c_{jx})^2+(y^i-c_{jy})^2}{2\sigma_j^2}\right)$:

\begin{table}[H]
\centering
\caption{Базисные функции, предсказания и невязки (it=0)}
\begin{tabular}{cccccc}
\toprule
$i$ & $z^i$ & $G_1^i$ & $G_2^i$ & $z_{\mathrm{RBF}}^i$ & $\varepsilon^i$ \\
\midrule
1 & 0.25 & 0.921708 & 0.055966 & $+0.086574$ & $-0.163426$ \\
2 & 3.02 & 0.921708 & 0.278023 & $+0.064368$ & $-2.955632$ \\
3 & 5.72 & 0.457342 & 0.755325 & $-0.029798$ & $-5.749798$ \\
4 & 1.51 & 0.132895 & 0.999301 & $-0.086641$ & $-1.596641$ \\
5 & 0.28 & 0.023387 & 0.759745 & $-0.073636$ & $-0.353636$ \\
\bottomrule
\end{tabular}
\end{table}

\[
  L^{(0)} = \frac{1}{2\cdot 5}\sum_{i=1}^5 (\varepsilon^i)^2 = 4.4497.
\]

\subsection{Аналитические выражения частных производных}

\begin{align}
  \frac{\partial L}{\partial w_0} &= \frac{1}{n}\sum_{i}\varepsilon^i = -2.163826, \\[4pt]
  \frac{\partial L}{\partial w_1} &= \frac{1}{n}\sum_{i}\varepsilon^i G_1^i = -1.144988, \\[4pt]
  \frac{\partial L}{\partial w_2} &= \frac{1}{n}\sum_{i}\varepsilon^i G_2^i = -1.407609, \\[4pt]
  \frac{\partial L}{\partial c_{1x}} &= \frac{1}{n}\sum_{i}\varepsilon^i w_1 G_1^i \frac{x^i-c_{1x}}{\sigma_1^2} = -0.036258, \\[4pt]
  \frac{\partial L}{\partial c_{1y}} &= \frac{1}{n}\sum_{i}\varepsilon^i w_1 G_1^i \frac{y^i-c_{1y}}{\sigma_1^2} = -0.037243, \\[4pt]
  \frac{\partial L}{\partial c_{2x}} &= \frac{1}{n}\sum_{i}\varepsilon^i w_2 G_2^i \frac{x^i-c_{2x}}{\sigma_2^2} = -0.031536, \\[4pt]
  \frac{\partial L}{\partial c_{2y}} &= \frac{1}{n}\sum_{i}\varepsilon^i w_2 G_2^i \frac{y^i-c_{2y}}{\sigma_2^2} = -0.035772, \\[4pt]
  \frac{\partial L}{\partial \sigma_1} &= \frac{1}{n}\sum_{i}\varepsilon^i w_1 G_1^i \frac{(x^i-c_{1x})^2+(y^i-c_{1y})^2}{\sigma_1^3} = -0.059933, \\[4pt]
  \frac{\partial L}{\partial \sigma_2} &= \frac{1}{n}\sum_{i}\varepsilon^i w_2 G_2^i \frac{(x^i-c_{2x})^2+(y^i-c_{2y})^2}{\sigma_2^3} = +0.051675.
\end{align}

\subsection{Обновление параметров ($\eta = 0.2$)}

\[
  \theta^{(1)} = \theta^{(0)} - \eta\,\nabla L(\theta^{(0)}).
\]

\begin{table}[H]
\centering
\caption{Параметры после первой итерации}
\begin{tabular}{cccc}
\toprule
Параметр & $\theta^{(0)}$ & $\nabla L$ & $\theta^{(1)}$ \\
\midrule
$c_{1x}$  & 1.5850  & $-0.036258$ & 1.592252 \\
$c_{1y}$  & 1.4000  & $-0.037243$ & 1.407449 \\
$c_{2x}$  & 4.0700  & $-0.031536$ & 4.076307 \\
$c_{2y}$  & 4.1033  & $-0.035772$ & 4.110488 \\
$\sigma_1$ & 1.8360 & $-0.059933$ & 1.847987 \\
$\sigma_2$ & 1.8360 & $+0.051675$ & 1.825665 \\
$w_0$     & 0.0000  & $-2.163826$ & 0.432765 \\
$w_1$     & 0.1000  & $-1.144988$ & 0.328998 \\
$w_2$     & $-0.1000$ & $-1.407609$ & 0.181522 \\
\bottomrule
\end{tabular}
\end{table}


%=============================================================================
\section{Задание 5. Аппроксимация RBF-сетью (программно, Стратегия 2)}
%=============================================================================

\subsection{Стратегия}

\begin{enumerate}
  \item \textbf{Без учителя}: K-Means (реализован с нуля, без sklearn) для инициализации центров.
  \item \textbf{С учителем}: градиентный спуск для обновления \emph{всех} параметров
        $(c_{1x}, c_{1y}, c_{2x}, c_{2y}, \sigma_1, \sigma_2, w_0, w_1, w_2)$,
        Loss~--- MSE, $\eta = 0.01$, 5000 итераций.
\end{enumerate}

\subsection{Шаг 1. K-Means: найденные центры}

Инициализация: случайно выбраны точки 2 и 5 как начальные центры.

\begin{table}[H]
\centering
\caption{Итерации K-Means}
\begin{tabular}{cccc}
\toprule
it & $c_1$ & $c_2$ & shift \\
\midrule
0 & (2.1500, 1.8800) & (4.9400, 5.1500) & 0.681724 \\
1 & (2.1033, 1.9633) & (4.5350, 4.6100) & 0.000000 \\
\bottomrule
\end{tabular}
\end{table}

\[
  \boxed{c_1 = (2.1033,\;1.9633), \quad c_2 = (4.5350,\;4.6100).}
\]

\subsection{Шаг 2. Начальное приближение}

\[
  \sigma_1 = \sigma_2 = \frac{|c_1 - c_2|}{2} = 1.7971, \quad
  w_0 = 0.0,\; w_1 = 0.1,\; w_2 = -0.1, \quad \eta = 0.01.
\]

\subsection{Шаг 3. Градиентный спуск: кривая обучения}

\begin{table}[H]
\centering
\caption{Динамика обучения (выборочно)}
\begin{tabular}{ccccc}
\toprule
Итерация & Loss & $c_1$ & $c_2$ & $(w_0, w_1, w_2)$ \\
\midrule
0    & 4.377178 & (2.104, 1.964) & (4.535, 4.610) & (0.022, 0.115, $-$0.090) \\
500  & 0.660305 & (2.897, 2.804) & (4.227, 4.265) & (0.994, 2.570, 0.477) \\
1000 & 0.214444 & (2.917, 2.831) & (4.049, 4.065) & (0.690, 3.964, 0.381) \\
1500 & 0.073383 & (2.908, 2.818) & (3.964, 3.970) & (0.431, 4.740, 0.337) \\
2000 & 0.025350 & (2.902, 2.811) & (3.918, 3.917) & (0.269, 5.190, 0.322) \\
2500 & 0.008801 & (2.899, 2.807) & (3.892, 3.887) & (0.171, 5.453, 0.316) \\
3000 & 0.003075 & (2.897, 2.805) & (3.876, 3.869) & (0.113, 5.608, 0.311) \\
3500 & 0.001087 & (2.896, 2.804) & (3.866, 3.858) & (0.080, 5.699, 0.306) \\
4000 & 0.000394 & (2.895, 2.804) & (3.860, 3.852) & (0.061, 5.753, 0.302) \\
4500 & 0.000150 & (2.895, 2.804) & (3.856, 3.848) & (0.050, 5.786, 0.298) \\
4999 & 0.000062 & (2.8945, 2.8044) & (3.8540, 3.8448) & (0.0444, 5.8047, 0.2939) \\
\bottomrule
\end{tabular}
\end{table}

\subsection{Финальные параметры}

\begin{table}[H]
\centering
\caption{Финальные параметры RBF-сети}
\begin{tabular}{cc}
\toprule
Параметр & Значение \\
\midrule
$c_{1x}$ & 2.8945 \\
$c_{1y}$ & 2.8044 \\
$\sigma_1$ & 0.9913 \\
$c_{2x}$ & 3.8540 \\
$c_{2y}$ & 3.8448 \\
$\sigma_2$ & 2.0997 \\
$w_0$ & 0.0444 \\
$w_1$ & 5.8047 \\
$w_2$ & 0.2939 \\
Финальный MSE & 0.00006207 \\
\bottomrule
\end{tabular}
\end{table}

\subsection{Аналитический вид модели}

\begin{equation}
  z_{\mathrm{RBF}}(x,y) = 0.0444
  + 5.8047\,\exp\!\left(-\frac{(x-2.8945)^2+(y-2.8044)^2}{2\cdot 0.9913^2}\right)
  + 0.2939\,\exp\!\left(-\frac{(x-3.8540)^2+(y-3.8448)^2}{2\cdot 2.0997^2}\right).
\end{equation}

\subsection{Невязки на каждом объекте}

\begin{table}[H]
\centering
\caption{Невязки RBF-сети}
\begin{tabular}{ccccc}
\toprule
$i$ & $(x^i, y^i)$ & $z^i$ & $z_{\mathrm{RBF}}^i$ & невязка \\
\midrule
1 & (1.02, 0.92) & 0.25 & 0.2488 & $-0.0012$ \\
2 & (2.15, 1.88) & 3.02 & 3.0156 & $-0.0044$ \\
3 & (3.14, 3.09) & 5.72 & 5.7051 & $-0.0149$ \\
4 & (4.13, 4.07) & 1.51 & 1.5161 & $+0.0061$ \\
5 & (4.94, 5.15) & 0.28 & 0.2984 & $+0.0184$ \\
\bottomrule
\end{tabular}
\end{table}

Невязки малы (порядка $10^{-2}$) --- RBF-сеть хорошо аппроксимирует данные за 5000 итераций.

\subsection{Визуализация}

\begin{figure}[H]
  \centering
  \includegraphics[width=\textwidth]{задание5_rbf.png}
  \caption{Задание 5: кривая обучения RBF-сети, 3D-поверхность и линии уровня}
\end{figure}



\end{document}
